\section{Связь с ВКР и дальнейшие шаги}

\textbf{Где здесь планирование.} Калибровка даёт гарантию, не зависящую от планировщика: любой путь, соблюдающий раздутые запретные зоны, безопасен с вероятностью не ниже $1-\alpha$. Но раздувание сужает свободное пространство и создаёт узкие проходы. Поэтому вклад планировщика измеряется не гарантией, а тем, как часто и как быстро он решает \emph{ужесточённую} задачу. Это постановка для ВКР по планированию движения в условиях неопределённости (научный руководитель И.~С.~Довгополик): двунаправленный RRT$^\ast$ с угловой динамической зоной выборки, стоимостью ненадёжности карты и откалиброванным запасом, который поставляет модуль из этой работы. Предварительный эксперимент раздела~\ref{sec:bridge} показывает, как решаемость за фиксированный бюджет зависит от масштаба запаса.

\textbf{Мост к SLAM (второй год).} Постановка~\eqref{eq:miss}--\eqref{eq:guarantee} не меняется при замене синтетического дрейфа реальной ошибкой оценщика позы; это уже проверено на визуальной одометрии. Следующие шаги:
\begin{enumerate}[nosep]
  \item Заменить одометрию кадр к кадру (разд.~\ref{sec:variants}) на SLAM с замыканиями циклов --- ORB-SLAM3, RTAB-Map или SLAM лаборатории. Это промежуточный режим ошибки позы, в котором гарантия может стать полезной: ошибка ограничена, но не нулевая.
  \item Увеличить число калибровочных сцен. Сейчас при 13 сценах запас определяет худшая сцена (разд.~\ref{sec:alpha}, \ref{sec:variants}). Данные OSMa-Bench++ (40 сцен, 4 условия освещения) позволят и работать при $\alpha = 0{,}05$, и проверить ось освещения.
  \item Сделать запас зависящим от объекта, а не одним на класс: нормированный счёт, у которого радиус растёт с неуверенностью найденного пятна. Диагностика показывает связь промаха велосипеда с уверенностью карты ($\rho \approx -0{,}5$).
  \item Заменить детектор рамок полным картографом (ConceptGraphs, BBQ) и проверить, сохраняются ли полные промахи и подмены класса целого объекта.
  \item Учесть ошибку локализации во время исполнения плана: сейчас гарантия относится к системе координат планировщика в момент запроса.
\end{enumerate}

\textbf{Открытый вопрос для лаборатории.} Если подмена класса целого объекта (тумба под ТВ $\to$ скамейка, ковёр) встречается чаще, чем в $\alpha$ доле сцен, никакая калибровка класс-специфичной области не даст полезной гарантии для этого класса: честный ответ --- отказ. Это перекладывает вопрос с планировщика на картограф: нужна не только высокая mIoU, но и редкость уверенных подмен для классов, важных для безопасности. Такую метрику (доля сцен с полной подменой класса) естественно добавить в протокол OSMa-Bench.
